\documentclass[a4paper]{article}
\usepackage[utf8]{inputenc}
\usepackage[russian,english]{babel}
\usepackage[T2A]{fontenc}
\usepackage[left=10mm, top=20mm, right=18mm, bottom=15mm, footskip=10mm]{geometry}
\usepackage{indentfirst}
\usepackage{amsmath,amssymb}
\usepackage[italicdiff]{physics}
\usepackage{graphicx}
\graphicspath{{images/}}
\DeclareGraphicsExtensions{.pdf,.png,.jpg}
\usepackage{wrapfig}
\usepackage{float}
\usepackage{subcaption}
\usepackage{enumitem}
\usepackage{caption}
\usepackage{array}
\captionsetup[figure]{name=Рисунок}
\captionsetup[table]{name=Таблица}

\title{\underline{Отчет о выполненой лабораторной работе 4.3.3}}
\author{Воронин Денис, Б04-407}

\begin{document}

\maketitle

\begin{center}
\textbf{\Large Исследование разрешающей способности
микроскопа методом Аббе}
\end{center}

\textbf{Цель работы:} определение дифракционного предела разрешения
объектива микроскопа методом Аббе.

\textbf{В работе используются:}  лазер; кассета с набором сеток разного
периода; линзы; щель с микрометрическим винтом; оптический стол
c набором рейтеров и крепёжных винтов; экран; линейка.

\textbf{Аннотация} должны будем определить периоды сеток сначала по их спектру на удаленном экране, затем по увеличенному с помощью модели микроскопа
изображению сеток на экране и по результатам измерения разрешающей способности микрскопа; наблюдать явления саморепродукции, пространственной фильтрации и мультиплицирования.

\section{Теоретические сведения}

Для иммерсионного микроскопа разрешающая способность объектива при
некогерентном освещении $$ \ell_{\min } \approx \frac{0.61 \lambda}{n \sin u}
$$ где $u$ -- апертурный угол объектива микроскопа (угол между оптической осью
и лучом, направленным из центра объекта в край линзы).

Метод Аббе для оценки разрешающей способности состоит в разделении хода лучей
на две части: сначала рассматривается картина в задней фокальной плоскости $F$
объектива; она называется первичным изображением или фурье-образом. Это
первичное изображение рассматривается как источник волн (согласно принципу
Гюйгенса-Френеля), создающий изображение в плоскости $P_{2}$, сопряжённой
плоскости предмета - вторичное изображение. Первичное изображение есть картина
дифракции Фраунгофера (на дифракционной решётке), если её период $d$, то для
направления максимальной интенсивности $\varphi_{m}$ $$ d \sin \varphi_{m}=m
    \lambda $$ При этом проходят пучки только с $\varphi_{m}<u .$ Можно условием
разрешения считать, что $u>\varphi_{1}$, иначе говоря $$ \sin u \geq \lambda /
    d $$ Или $$ d \geq \frac{\lambda}{\sin u} \approx \frac{\lambda}{D / 2 f} $$
где $D-$ диаметр линзы, $f-$ фокусное расстояние. Двумерную решётку можно
рассматривать как две перпендикулярные друг другу, для максимумов которых
выполняется соотношение $$ d \sin \varphi_{x}=m_{x} \lambda, \quad d \sin
    \varphi_{y}=m_{y} \lambda $$

\newpage

\section{Эксперементальная установка}

\begin{figure}[!hb]
	\centering
	\includegraphics[width=0.9\linewidth]{p1.png}
	% \caption{Микрометирический винт}
	\label{pic_3}
\end{figure}

Схема установки приведена на рисунке. Предметом $P_{1}$ служат
сетки в кассете $C$. Линза Л1 -- длиннофокусная, а Л2 -- короткофокусная. В F
Устанавливаются диафрагмы D, с помощью сеток с разными периодами $d$ и щелевой
диафрагмы можно проверить третье соотношение. Период сеток может быть измерен
либо по расстоянию между дифракционными максимумами на экране, либо по
увеличенному с помощью микроскопа изображению сетки на экране. Пространственную
фильтрацию (получение наклонного изображение решётки) можно получить с помощью
подбора угла наклона и ширины вспомогательной щели.

\section{Ход работы}

\subsection{Определение периода решеток по их пространственному спектру}

\textbf{1.} Включили лазер в сеть. Выделили самый яркий луч и следили за ним.

\textbf{2.} Закрепили касету с двумерными решетками близи выходного окна лазера. Пронаблюдали дифракционные картинки для разных сеток.
Для определения расстояния между соседними дифракционными максимумами измерили расстояние между удаленными друг от друга максимумами 
(вертикальными) и число промежутков между ними:

\begin{tabular}{|c|c|c|}
\hline
Номер решетки & $\uparrow $ & $\downarrow $ \\
\hline
1    & 80 мм (1 пром.)   & 80 мм (1 пром)   \\
\hline
2    &   33 мм (1 пром)  & 33 мм (1 пром)    \\
\hline
    & 63 мм (2 пром)    & 64 мм (2 пром)    \\
\hline
    & 94 мм (3 пром)    & 94 мм (3 пром)    \\
\hline
3    & 16 мм (1 пром)    & 16 мм (1 пром)    \\
\hline
    & 32 мм (2 пром)    & 31 мм (2 пром)    \\
\hline
    & 47 мм (3 пром)    & 47 мм (3 пром)    \\
\hline
    & 62 мм (4 пром)    & 62 мм (4 пром)    \\
\hline
    & 76 мм (5 пром)    & 77 мм (5 пром)    \\
\hline
    & 94 мм (6 пром)    & 93 мм (6 пром)    \\
\hline
\end{tabular}

\par

\textbf{3.} Расстояние от экрана до сетки $ H = 150 $ см.\par

Длина волны $\lambda_{\text{зел}} = 532 $ нм.

\[d = \frac{m*\lambda}{sin\theta}\]

Для первой сетки: d = 13 $\pm 0.3$ мкм \par
Для второй сетки: d = 25.6 $\pm 0.5$ мкм \par
Для первой сетки: d = 50 $\pm 0.9$ мкм \par

\subsection{Определение периода решеток по изображению, увеличенному с помощью модели микроскопа}

\textbf{4.} Собрали модель проекционного микроскопа. Отцентрировали систему. 
Получили резкое изображение $h\nu $: 

\begin{figure}[H]
	\center{\includegraphics[scale=0.25]{p1.jpg}}
  \caption{Настройка микроскопа}
	\label{fig:image1}
\end{figure}



\textbf{5.} Определили расстояния $a_{1} = 13.5$ см, $a_{2} = 2.5$ см, $b_{1} = 67$ см, $b_{2} = 64.5$ см. \par
Увеличение системы: 

\[\text{Г} = \frac{b_{1}b_{2}}{a_{1}a_{2}} = \frac{67*64.5}{13.5*2.5} = 128\]

Периоды изображений сеток на экране:

Для первой сетки: d = 13.0 $\pm 0.3$ мкм \par
Для второй сетки: d = 25.0 $\pm 0.5$ мкм \par
Для первой сетки: d = 43 $\pm 0.8$ мкм \par

\textbf{Диаметр волоса:} \par

\[tg\theta = \frac{\lambda}{d} \Rightarrow d = \frac{\lambda}{tg\theta} = \frac{532* 10^{-9}}{\frac{0.01}{1.5}} = 79 \text{мкм}\]

\section{Определение периодов решетки по оценке разрешающей способности микроскопа}
\textbf{6.} Поместили щелевую диафрагму с микрометрическим винтом в фокальной плоскости F линзы $\text{Л}_{1}$. Определили для каждой решетки минимальный размер диафрагмы D, при
котором на экране еще видно изображение сетки:\par

\textbf{Открытие щели}: 0.1 мм \par
\textbf{Горизонтальное направление}: 1.24 мм \par
\textbf{Вертикальное направление}: 3.07 мм \par

Для второй и третьей щели размера щели не хватает, поэтому не получилось измерить.

\subsection{Пространственная фильтрация и мультиплицирование}

Проделали качественный опыт по пространственной фильтрации.

\begin{figure}[H]
	\center{\includegraphics[scale=0.2]{p2.jpg}}
  \caption{Качественный опыт}
	\label{fig:image1}
\end{figure}

\begin{figure}[H]
	\center{\includegraphics[scale=0.2]{p3.jpg}}
  \caption{Качественный опыт}
	\label{fig:image1}
\end{figure}

\begin{figure}[H]
	\center{\includegraphics[scale=0.2]{p4.jpg}}
  \caption{Качественный опыт}
	\label{fig:image1}
\end{figure}

\section{Обработка результатов}

Рассчитаем минимально расстояние d, разрешаемое микроскопом:

\[d \approx \frac{\lambda*2f}{D} = \frac{532*10^{9}*2*0.135}{2.93*10^{-3}} = 49.0 \pm 0.7 \text{мкм} \]

Проверим теорию Аббе, построив график зависимости $d = f(1/D)$ взяв периоды сеток, определенных по спектру:

\begin{table}[h]
    \centering
    \begin{tabular}{|l|l|l|}
        \hline
        № реш. & $D$, мм & $d_{\text{эксп}}$, мкм \\
        \hline
        1 & $4$ & $13$ \\
        \hline
        2 & 2,20 & 25 \\
        \hline
        3 & 1,24 & 50 \\
        \hline
    \end{tabular}
\end{table}

\begin{figure}[H]
	\center{\includegraphics[scale=0.5]{F1.png}}
  \caption{График зависимости $d = f(1/D)$}
	\label{fig:image1}
\end{figure}

Если сравнить коэффициент наклона у графика, и посчитать значение по формуле $\lambda*f$, то получим близкие результаты, отличающиеся на 0.02 \%.

\section{Вывод} 
Определил дифракционный предел разрешения объектива микроскопа методом Аббе. Измерил периоды сеток сначала по их спектру на удаленном экране, затем по увеличенному с помощью модели микроскопа изображению сеток.
Пронаблюдал явление саморепродукции, пространственной фильтрации и мультиплицирования.

\end{document}